\chapter{Auditoría del desarrollo}
\label{anexo:desarrollo}

Este anexo conserva sólo los defectos que pudieron cambiar la validez o la interpretación del
estudio. No es una cronología del código ni una lista de incidencias menores. Su propósito es
mostrar qué tipos de error eran capaces de producir resultados plausibles, cómo se detectaron y qué
contrato quedó después para impedir su regreso. La historia completa permanece en la bitácora del
repositorio.

{\footnotesize
\input{tables/tD_defectos_validez.tex}
}

\section{Cuatro defectos que producían resultados plausibles}

\paragraph{Una regla estadística correcta en apariencia.} La puerta de no inferioridad construía su
intervalo sobre una transformación que ensanchaba el castigo conforme aumentaba la ventaja del
candidato, de modo que ningún retador podía resultar elegible ni siquiera dominando al incumbente.
No fallaba nada: los runs terminaban y las métricas parecían razonables. El efecto material se vio
en el conjunto de variables, donde la versión defectuosa retenía el reducido con Rank-IC 0,0730 y,
corregida la comparación, el amplio alcanzó 0,0958. La corrección fue reescribir la puerta para que
el objeto estadístico fuese la serie de diferencias emparejadas por cohorte, y de ahí salieron tres
contratos: el signo de una ventaja se conserva al construir su intervalo, un candidato estrictamente
mejor no puede volverse menos elegible porque crezca su incertidumbre, y toda decisión registra
tamaño de muestra, fracción de victorias y límites del intervalo.

\paragraph{Configuración declarada que no cambia la ciencia.} Una opción puede aparecer en el
catálogo, viajar por la API y quedar guardada sin alterar los datos que consume el modelo. Ocurrió
con una ablación: seis factores de precio llegaban al agente de tendencia fuera del condicional que
debía retirarlos, de modo que la interfaz decía haber eliminado una familia de información que el
modelo seguía recibiendo. El riesgo no es la etiqueta incorrecta sino que una conclusión del tipo
«la información de precio no importa» resultaría irrefutable desde las métricas, porque los dos
modelos serían casi el mismo. La misma lógica detectó que un parámetro de submuestreo no activaba el
\textit{bagging} por faltarle su frecuencia. La lección: auditar el valor declarado prueba la
interfaz; auditar el consumidor prueba el experimento.

\paragraph{Identidades y ventanas que mezclaban evidencias.} La primera clave de evaluación no
incorporaba la identidad del panel, así que dos conjuntos de datos distintos podían compartirla; el
indicio fue encontrar bajo la misma identidad un CAGR de 0,1468 y otro de 0,1692. En paralelo, las
métricas económicas se calculaban sobre la curva completa y mezclaban selección con la era
reservada, lo que impedía saber qué parte era evidencia disponible durante el desarrollo. La salida
se segmentó en selección, reserva y curva completa. También se corrigió el exceso geométrico, que
restaba dos CAGR en vez de comparar los factores de riqueza:
\[
 \alpha_{\mathrm{geom}}=
 \left(\frac{W_T^{\mathrm{cartera}}}{W_T^{\mathrm{SPY}}}\right)^{1/Y}-1.
\]

\paragraph{Una simulación que termina no es una simulación válida.} Los defectos de cartera más
peligrosos no lanzaban excepciones: una posición sin precio podía marcarse plana, una venta podía
pagar costes sin ejecutarse, una señal ausente podía bloquear una rotación. Cada caso produce una
curva suave y aparentemente conservadora que altera el contrafactual. La convención final penaliza
una desaparición sin precio con un $-30$\,\% y liquida contra efectivo, sólo cobra comisión si hubo
orden ejecutada, y registra en cada venta la alternativa a la que se reasigna el capital. La prueba
de carteras aleatorias mostró lo mismo: el nulo inicial llegaba a un percentil 95 de CAGR del
107\,\% porque aceptaba historias incompletas y no pagaba las mismas fricciones, de modo que no era
evidencia contra el modelo sino un contrafactual distinto.

\section{Matriz de contratos que quedó en el sistema}

\begin{table}[htbp]
  \centering
  \small
  \begin{tabularx}{\textwidth}{p{3.0cm}X X}
    \toprule
    Frontera & Invariante & Fallo que evita \\
    \midrule
    Catálogo--API & Rechazo de claves y combinaciones desconocidas & Hipótesis no declaradas. \\
    Configuración--datos & Columnas efectivas iguales a las autorizadas & Ablaciones ficticias. \\
    Datos--label & Publicación y cierre anteriores al entrenamiento & Fuga temporal. \\
    Evaluación--caché & Dataset, catálogo y clave coincidentes & Reutilización entre paneles. \\
    Selección--reserva & Ningún artefacto de rejilla contiene 2025--2026 & Mirada anticipada. \\
    Orden--contabilidad & Patrimonio reconstruible desde operaciones & Costes o ventas fantasma. \\
    Evidencia--LaTeX & Activo y fuente declarados en el manifiesto & Mezcla de estudios. \\
    \bottomrule
  \end{tabularx}
  \caption{Contratos de validación derivados de los defectos materiales.}
  \label{tab:contratos-validacion}
\end{table}

La Tabla~\ref{tab:contratos-validacion} recoge los contratos que quedaron en el sistema tras esas
correcciones. Ninguna de ellas aumenta el Rank-IC ni es evidencia económica. Su valor está en delimitar qué
puede creerse del resultado: el modelo recibió las variables que declaraba, las filas pertenecían al
tiempo que les correspondía, la búsqueda comparó candidatos bajo una regla coherente y la cartera
pagó sólo operaciones que ocurrieron. La conclusión de fondo es que la plausibilidad no sirve como
prueba: un número financiero puede parecer razonable y estar midiendo otro universo, otra ventana o
una hipótesis que nunca llegó a activarse.
